\icode{E-MI-SI-L3-261029-001}
如图所示，\(A\)、\(B\)两物体的质量分别为\(m_A\)、\(m_B\)，且\(m_A > m_B\)，整个系统处于静止状态，滑轮的质量和一切摩擦均不计，如果绳一端由\(Q\)点缓慢地向左移到\(P\)点，\(PQ\)距离为\(L\)，整个系统重新平衡后，物体的高度和两滑轮间的绳与水平方向的夹角\(\theta\)变化情况是\paren[]

\begin{tikzpicture}[font=\scriptsize]
    \draw ($(-60:.1)+(.3,0)$) arc(0:30:.3);
    \draw[dashed,cyan] (-60:.1)--++(1,0);
	\draw[magenta] (-120:.1)--++(150:3)coordinate(Q) (-60:.1)--++(30:2.6)coordinate(A)--++(-60:.1)coordinate(B)--++(0:.1)coordinate(C)--++(0,-1.5)coordinate(D);
	\draw[fill=white] (0,0) circle(.1) (B) circle(.1);
	\foreach \x in {-3.6,-3.4,...,3}
		\draw[gray] (\x,{1.5-0.05*sqrt(3)})--++(.2,.2);
	\draw[thick] (-3.6,{1.5-0.05*sqrt(3)})--++(6.8,0) (B)--(B|-Q) (0,0)--(0,-.6);
	\coordinate (E)at($(D)+(-.2,0)$);
	\coordinate (F)at($(D)+(.2,-.4)$);
	\draw[fill=orange!60] (-.15,-.6) rectangle(.15,-.9)node[midway,black]{\(B\)} (E) rectangle(F)node[midway,black]{\(A\)};
	\fill (0,0) circle(.05) (B) circle(.05) (Q)node[below]{\(Q\)} circle(.05) ($(Q)+(-.8,0)$)node[below]{\(P\)} circle(.05);
	\node at(.6,.05){\(\theta\)};
\end{tikzpicture}
\begin{choices}
	\item 物体\(A\)上升，上升的高度为\(\frac{L}{\cos\theta}\)，\(\theta\)角不变 
	\item 物体\(A\)下降，下降的高度为\(\frac{L}{2\cos\theta}\)，\(\theta\)角不变
	\item 物体\(A\)上升，上升的高度为\(\frac{2L}{\cos\theta}\)，\(\theta\)角不变
	\item 物体\(A\)的高度不变，\(\theta\)角变小
\end{choices}
\begin{answer}
	A\qquad 提示：受力分析可知：\(2 m_Ag\sin\theta = m_Bg\)，新建立的平衡位置如图所示

\begin{tikzpicture}[font=\scriptsize]
    \draw ($(-60:.1)+(.3,0)$) arc(0:30:.3);
    \draw[dashed,cyan] (-60:.1)--++(1,0);
	\draw[dashed,magenta] (-120:.1)--++(150:3)coordinate(Q) (-120:.1)--++({-sqrt(3)},0);
    \draw[magenta] (-60:.1)--++(30:2.6)coordinate(A)--++(-60:.1)coordinate(B)--++(0:.1)coordinate(C)--++(0,-.5)coordinate(D);
    \draw[magenta] ($(Q)+({-sqrt(3)},0)$)coordinate(P)--++(-30:4)coordinate(p) ($(p)+(.1,0)$)--(-60:.1);
    \draw ($(p)+(60:.1)$) circle(.1) ($(p)+(60:.1)$)--++(0,-.6)coordinate(b);
    \draw[fill=orange!60] ($(b)+(-.15,0)$) rectangle ($(b)+(.15,-.3)$)node[midway,black]{\(B\)};
    \fill[gray] ($(p)+(60:.1)$) circle(.05);
	\draw[fill=white] (B) circle(.1);
	\foreach \x in {-4.8,-4.6,...,3}
		\draw[gray] (\x,{1.5-0.05*sqrt(3)})--++(.2,.2);
	\draw[thick] (-4.8,{1.5-0.05*sqrt(3)})--++(8,0) (B)--(B|-Q);
    \draw[gray] (0,0)--(0,-.6) (0,0) circle(.1) (0,0) circle(.05);
	\coordinate (E)at($(D)+(-.2,0)$);
	\coordinate (F)at($(D)+(.2,-.4)$);
	\draw[fill=orange!30,draw=gray] (-.15,-.6) rectangle(.15,-.9);
    \draw[fill=orange!60] (E) rectangle(F)node[midway,black]{\(A\)};
	\fill (B) circle(.05) (Q)node[below]{\(Q\)} circle(.05) (P)node[below]{\(P\)} circle(.05);
	\node at(.6,.05){\(\theta\)};
\end{tikzpicture}
\end{answer}